\documentclass[11pt,a4paper]{article}
\usepackage[margin=2.2cm]{geometry}
\usepackage{booktabs,enumitem,graphicx}
\usepackage[T1]{fontenc}
\setlength{\parskip}{0.5em}\setlength{\parindent}{0pt}
\title{Kinect Lite: Explanation of Results}
\date{}
\begin{document}
\maketitle

\section*{What was tested}
Kinect depth images contain \textbf{spikes} (pixels with a wildly wrong distance) and \textbf{holes} (pixels with no distance). Two cleaners were compared: our \textbf{9$\times$9 median} and the \textbf{Fu et al.\ (2012) paper pipeline}. Each cleaned sequence was fed to two trackers, \textbf{depth-only ICP} and \textbf{RGB-D odometry} (colour + depth), and the estimated camera path was compared with the true path. ATE is total drift (cm); RPE is per-step error (mm); lower is better.

\section*{1. Depth accuracy (ICL, against true depth)}
\begin{center}
\begin{tabular}{lccc}\toprule
Method & MAE (mm) & RMSE (mm) & within 1 cm\\midrule
No cleaning & 39.2 & 187.5 & 42.2\%\
Median 9$\times$9 (ours) & \textbf{18.8} & \textbf{55.4} & \textbf{44.7\%}\
Paper (Fu et al.) & 37.7 & 186.0 & 33.7\%\\bottomrule
\end{tabular}
\end{center}
\begin{itemize}[nosep]
\item The raw RMSE (187 mm) is almost 5$\times$ its MAE, so \textbf{spikes dominate the error}.
\item The median votes spikes out: the error roughly halves and RMSE drops to 55 mm.
\item The paper's bilateral filter weights neighbours by depth similarity, so a spike only averages with itself and survives. Its error barely changes.
\end{itemize}

\section*{2. Camera tracking (ATE in cm)}
\begin{center}
\begin{tabular}{lcccc}\toprule
Method & ICL ICP & ICL RGB-D & TUM ICP & TUM RGB-D\\midrule
True depth & 0.19 & 0.81 & -- & --\
No cleaning & 17.07 & 23.59 & 13.97 & 2.14\
Median (ours) & 14.94 & \textbf{0.98} & 13.32 & \textbf{2.08}\
Paper & \textbf{0.95} & 26.0 & \textbf{3.31} & 2.10\\bottomrule
\end{tabular}
\end{center}
The best cleaner depends on the tracker:
\begin{itemize}[nosep]
\item \textbf{Depth-only ICP favours the paper.} ICP discards far-off matches, so leftover spikes cost little, but it relies on surface normals. The paper's smooth surfaces give good normals; the median leaves a striped error pattern that biases each step slightly and accumulates as drift.
\item \textbf{RGB-D odometry favours the median.} It projects every pixel using its depth, so each surviving spike drags the alignment. With the paper's output 49 of 149 steps were off by more than 20 mm; with the median, none. The median's 0.98 cm is close to the 0.81 cm from perfect depth.
\item \textbf{TUM RGB-D is a tie} (2.08--2.14 cm): real sensor noise has fewer spikes than ICL's simulated noise, so cleaning matters less.
\item The paper's ambiguous hole-prediction step does not matter: all three readings agree within 0.06 cm, since holes are only about 1\% of pixels.
\end{itemize}

\section*{3. Should holes be filled?}
\begin{center}
\begin{tabular}{lcccc}\toprule
Fill method & Fill err (mm) & ICL RGB-D & TUM RGB-D & Time/frame\\midrule
None (keep holes) & -- & \textbf{0.98} & \textbf{2.08} & 0 ms\
Telea inpainting & 87.0 & 1.44 & 3.13 & 210 ms\
Navier--Stokes & 22.5 & 1.46 & 3.89 & 224 ms\
Nearest pixel & 22.3 & 1.43 & 4.16 & 50 ms\
Region growing (small) & 15.5 & 1.01 & 2.23 & 52 ms\
Temporal (paper) & \textbf{12.2} & 1.24 & 2.51 & 7.9 s\\bottomrule
\end{tabular}
\end{center}
No fill improves RGB-D odometry. A filled pixel is a guess, and the tracker does better ignoring a missing pixel than using a wrong one. Telea, the obvious choice, is among the least accurate because it blends the near and far sides of a depth edge. The fills that change the least hurt the least.

\section*{Why the best cleaner flips with the tracker}
The depth image has two kinds of mistakes:
\begin{itemize}[nosep]
\item \textbf{Spikes:} a few pixels with a very wrong distance.
\item \textbf{Roughness:} surfaces are slightly bumpy instead of flat, even where the wall is flat.
\end{itemize}
Each cleaner fixes only one of them: the \textbf{median removes spikes} but leaves the roughness; the \textbf{paper smooths surfaces} but leaves the spikes.

Each tracker is bothered by only one of them: \textbf{depth-only ICP} lines up surfaces, so it cares about roughness and mostly ignores spikes; \textbf{RGB-D odometry} uses depth at every pixel, so it cares about spikes, while the colour image makes it less sensitive to roughness.
\begin{center}
\begin{tabular}{lll}\toprule
 & Median (fixes spikes) & Paper (fixes roughness)\\\midrule
Depth-only ICP (bothered by roughness) & still drifts: 14.94 cm & works: 0.95 cm\\
RGB-D odometry (bothered by spikes) & works: 0.98 cm & fails: 26.0 cm\\\bottomrule
\end{tabular}
\end{center}
So the median fixes the problem RGB-D has, and the paper fixes the problem ICP has. \emph{Caveat:} the ``roughness hurts ICP'' part is the project's explanation, not directly proven by the experiments. The numbers only show that the paper helps ICP a lot.

\section*{Conclusion}
\begin{center}
\begin{tabular}{p{4.2cm}p{4cm}p{6cm}}\toprule
System uses & Use & Why\\midrule
RGB-D odometry & 9$\times$9 median, no hole filling & Error halved, best tracking, under 1 min for 450 frames\
Depth-only ICP & Paper pipeline & 4--16$\times$ less drift, but about 26 min for 450 frames\
Display only & Median + Navier--Stokes or nearest & About 4$\times$ more accurate than Telea; do not feed to the tracker\\bottomrule
\end{tabular}
\end{center}
\end{document}
